%% chapter5.tex: Conclusions and Future Work

\chapter{Conclusions and Future Work}
\label{chap:conclusion}

\section{Conclusions}

This dissertation presented the design, implementation, and experimental
evaluation of a routing mechanism for a mobile ad-hoc network built on top
of the RA-TDMAs+ TDMA protocol. The work was motivated by the need to provide
reliable multi-hop IP communication for a teleoperated mobile robot streaming
video to a base station in environments without fixed network infrastructure.

The main technical contributions were: (1) a link quality metric derived from
TDMA beacon reception statistics; (2) an MST-based routing algorithm that
derives per-destination next-hop decisions from the quality-weighted topology
graph; (3) a dual routing table architecture that uses the Linux kernel's
native IP forwarding for relay traffic, eliminating the need for a user-space
relay plane; and (4) a video quality feedback loop that enables relay
activation independently of the TDMA topology convergence cycle.

Experimental evaluation on physical hardware demonstrated successful relay
activation after link failure in 67\% of the tested cases, with a mean
convergence time of 2.76\,s. The dominant bottleneck was identified as the
video loss detection timeout at the base station (1.5\,s), not the routing
computation or kernel route installation, which are completed in microseconds.
Under relay conditions, round-trip latency approximately doubled (from 48\,ms
to 97\,ms) as expected from the two-hop path, and the video stream delivery
rhythm was fully preserved without additional jitter at the frame level.

The results confirm the viability of kernel-level IP routing as an alternative
to data-link-layer forwarding for RA-TDMAs+ networks. The approach provides
standard-tool observability (tcpdump, traceroute, netstat), compatibility
with any IP protocol, and integration with the Linux networking stack's
production features, while meeting the timing constraints imposed by the
50\,ms TDMA slot.

\section{Future Work}
\label{sec:future}

Several directions for future work are identified:

\begin{itemize}

  \item \textbf{Reduced convergence time}: The 1.5\,s video loss timeout
  was chosen conservatively. Adaptive timeout mechanisms that distinguish
  between transient packet loss and link failure could reduce the detection
  time to 300--500\,ms without increasing the false-positive rate.

  \item \textbf{Proactive relay pre-positioning}: When the link quality metric
  for the direct link falls below a configurable threshold (but has not yet
  reached zero), the system could pre-install relay routes in parallel with
  the direct route, reducing convergence to a single TDMA frame.

  \item \textbf{Larger networks}: The current implementation was validated
  with three nodes. Extension to larger networks requires addressing the MST
  recomputation overhead (currently $O(N^2)$, negligible for $N \leq 10$) and
  the TCP connection management overhead ($O(N^2)$ connections per node).

  \item \textbf{Multi-path relay}: When multiple relay nodes are available,
  the system currently uses a single relay path. Load balancing across
  multiple paths could increase throughput and resilience.

  \item \textbf{Energy-aware routing}: Incorporating battery state into the
  link quality metric would allow the MST computation to prefer paths through
  nodes with higher remaining energy, extending network lifetime in
  battery-constrained deployments.

\end{itemize}
